%! crea ambiente "mybox" para cajas de texto
\newtcolorbox{mybox}
{
  colframe = black!40,
  colback  = black!3
}
\newtcolorbox{mybox2}
{
  colframe = black!40,
  colback  = white
}


%! redefinición de ambientes
\newtranslation[to=spanish]{Section}{Secci\'on}
\newtranslation[to=spanish]{Subsection}{Subsecci\'on}

%! ambientes theorem
\newtheorem{thm}{Teorema}
\newtheorem{cor}[thm]{Corolario}
\newtheorem{lem}[thm]{Lema}
\newtheorem{propo}[thm]{Proposición}
\newtheorem{prop}[thm]{Propiedades} %ambiente "proposition"
\theoremstyle{definition}
\newtheorem{defn}[thm]{Definición}
\newtheorem{eg}{Ejemplo}[section]
\newtheorem{egs}{Ejemplos}[section]
\theoremstyle{remark}
\newtheorem{rem}[thm]{Nota}
\numberwithin{equation}{section}

%! funciones matemáticas
\newcommand{\sen}{\operatorname{sen}}
\newcommand{\tg}{\operatorname{tg}}
\newcommand{\ctg}{\operatorname{ctg}}
\newcommand{\sech}{\operatorname{sech}}
\newcommand{\csch}{\operatorname{csch}}
\newcommand{\ctgh}{\operatorname{ctgh}}
\newcommand{\cosec}{\operatorname{cosec}}
\newcommand{\Arctg}{\operatorname{Arctg}}
\newcommand{\Arccot}{\operatorname{Arccot}}
\newcommand{\Arccos}{\operatorname{Arccos}}
\newcommand{\Arcsen}{\operatorname{Arcsen}}
\newcommand{\To}{\longrightarrow}
\newcommand{\ssi}{\leftrightarrow}
\newcommand{\Ssi}{\Leftrightarrow}
\newcommand{\Real}{\mathds{R}}
\newcommand{\Ent}{\mathds{Z}}
\newcommand{\Nat}{\mathds{N}}
\newcommand{\Comp}{\mathds{C}}
\newcommand{\abs}[1]{\left|#1\right|}
\newcommand{\norm}[1]{\left\Vert#1\right\Vert}
\newcommand{\set}[1]{\left\{#1\right\}}
\newcommand{\gen}[1]{\left\langle #1\right\rangle}
\newcommand{\interior}[1]{\accentset{\circ}{#1}}
\newcommand{\prodin}[2]{\left\langle #1, #2\right\rangle}

%! intervalos, evita que los paréntesis se descuadren en la IDE
\newcommand{\intervaloAA}[1]{\left] #1 \right[}
\newcommand{\intervaloAC}[1]{\left] #1 \right]}
\newcommand{\intervaloCA}[1]{\left[ #1 \right[}
\newcommand{\intervaloCC}[1]{\left[ #1 \right]}

%! definición de operadores
\newcommand{\Ker}{\operatorname{Ker}}
\newcommand{\Ima}{\operatorname{Im}}

%! K para campos
\newcommand{\K}{\mathds{K}}

%! definición colores de texto
\newcommand{\red}[1]{{\color{red}{#1}}}
\newcommand{\blue}[1]{{\color{blue}{#1}}}
\definecolor{forestgreen}{rgb}{0.13, 0.55, 0.13}
\newcommand{\green}[1]{{\color{forestgreen}{#1}}}
\newcommand{\redbf}[1]{{\color{red}{\textbf{#1}}}}
\newcommand{\violet}[1]{{\color{violet}{#1}}}
\definecolor{col1}{RGB}{50,40,180}
